\section{通用视觉表示：从人类少样本归纳到可测目标}

在进入模型之前，必须先定义目标。一个表示若只能在预训练标签空间内工作，就更像专用特征；所谓 \term{general visual representation}（通用视觉表示），应让新任务在很少样本、较小适配成本下获得可靠性能，并覆盖自然图像、结构化场景、遥感、计数与几何等不同视觉问题。

\subsection{目标不是“看见”，而是快速适配}

讲者先提出一个面向真实使用者的愿景：未来普通人不需要编程，也能用少量示例教会系统处理新的视觉任务。这里的关键不是机器人外形，而是 adaptation interface——用户提供示例后，系统能否利用已有表示快速理解新分类规则。

\lecturefigure{slide-02-goal-general-representation.jpg}{研究目标是 general visual representation，而不是单个 benchmark 的最高分。}{官方视频 00:36。}

\begin{knowledgebox}{读图：representation 是任务前的公共资产}
预训练模型可写为特征映射 $h=f_{\theta}(x)$，下游任务再学习 $g_{\phi}(h)$。若 $f_{\theta}$ 真正通用，新的 $g_{\phi}$ 应只需要少量数据或较小参数更新。图中没有指定分类、检测或分割，正是在强调表示先于具体任务；衡量它必须看一组下游任务，而不是只看 upstream accuracy。
\end{knowledgebox}

\lecturefigure{slide-03-why-visual-representation.jpg}{好的视觉表示应“kick-start”大量含视觉输入的任务，并降低非专家的教学成本。}{官方视频 01:40。}

\begin{importantbox}{什么叫 kick-start}
Kick-start 不是零成本完成任务，而是显著降低三类成本：\textbf{sample cost}（需要多少标注示例）、\textbf{optimization cost}（需要训练多久、更新多少参数）和 \textbf{specification cost}（用户需要多精确地描述任务）。一个表示可能提高最终精度，却没有降低这些成本，因此仍不算易适配。
\end{importantbox}

\teachervoice{Lucas 用“孩子或父母能教会一个简单机器人”解释研究愿景，同时明确说视觉表示只是完整智能系统的一部分。课堂提示：把表示学习的价值写成用户适配成本，比泛称“模型更智能”更可检验。}

\subsection{人类视觉表示提供了哪种参照}

人类并非只擅长识别熟悉物体。讲者通过花朵、卫星图与抽象物体展示：即使类别语义陌生，人也能从少数示例推断规则。真正值得模仿的是从表示中组合出新判别边界的能力，而不是记住海量类别名称。读下面三组图时，应分别判断任务依赖长期语义经验、低层视觉统计，还是数量与关系等抽象规则；只有跨这些机制都能快速适配，才接近“通用”。

\lecturefigure{slide-04-human-visual-representation.jpg}{人类视觉表示被用作少样本归纳能力的参照。}{官方视频 01:48。}

\begin{knowledgebox}{读图：少样本成功可能来自不同机制}
花朵任务可利用长期视觉经验；卫星图需要迁移纹理与空间结构；抽象物体任务则可能依赖计数或关系规则。三者都表现为 few-shot classification，但需要的 invariance 不同。一个 benchmark 若只包含相似自然物体，无法区分“记住语义模板”和“形成通用视觉结构”。
\end{knowledgebox}

\lecturefigure{slide-05-vtab-flowers.jpg}{VTAB 示例中的花朵类别接近自然图像分布，少量样本即可形成类别概念。}{官方视频 02:00；Zhai et al. 2019。}

\begin{knowledgebox}{读图：先判断任务与预训练分布的距离}
花朵图片有真实纹理、颜色与轮廓，通常与 ImageNet 类预训练较接近。若模型在此类任务上表现好，可能来自可迁移表示，也可能来自上游类别/图像的语义重叠。读 benchmark 时第一步不是看分数，而是问 downstream 与 upstream 在图像来源、标签语义和视觉统计上有多近。
\end{knowledgebox}

\lecturefigure{slide-06-vtab-objects.jpg}{抽象物体任务要求根据数量、形状或关系归纳类别规则。}{官方视频 03:04；VTAB structured tasks。}

\begin{warningbox}{Few-shot 不是“只训练五张图”这么简单}
模型在下游只看到少量标注，并不代表它没有见过相似视觉模式。应区分 pretraining exposure、downstream labeled examples、hyperparameter tuning data 和 test examples。若在多个下游任务上反复调 recipe，整体系统使用的信息远多于单任务的 $k$ 个样本。
\end{warningbox}

\subsection{本章小结}

通用表示的核心是快速适配多种规则，而非单一分类精度。人类示例提示我们要同时覆盖熟悉语义、domain shift 与结构化推理；这直接引出 VTAB 的多任务测量框架。

